\documentclass[10pt,a4paper]{amsart}
\usepackage[margin=19mm]{geometry}
\usepackage{amsmath,amssymb,mathtools}
\usepackage[scr=rsfs,cal=euler]{mathalfa}
\usepackage{xfrac}
\usepackage[unicode,hidelinks,bookmarks=false]{hyperref}
\newcommand{\ie}{i.e.\ }
\newcommand{\cf}{cf.\ }
\newcommand{\resp}{respectively\ }
\newcommand{\Subsection}{\subsection}
\providecommand{\nobreakdash}{\nobreak\hbox{-}}
\setlength{\emergencystretch}{2em}
\multlinegap0pt
\usepackage{amssymb}
\usepackage{mathtools}
\usepackage[shortlabels]{enumitem}
\usepackage{float}
\usepackage{xcolor}
\usepackage{tikz}
\newtheorem{lemma}{Lemma}
\newtheorem{proposition}{Proposition}
\newtheorem{theorem}{Theorem}
\newcommand{\V}{\mathbf V}
\title{\textbf{The Prime Clockwork: A Dynamic Representation of Modular and Multiplicative Arithmetic}}
\author{Michael T.M. Emmerich}
\date{Source: arXiv:2609.03896v1 (CC BY 4.0)}
\begin{document}
\maketitle

\noindent\emph{Source and licence.} \textbf{The Prime Clockwork: A Dynamic Representation of Modular and Multiplicative Arithmetic}. Version arXiv:2609.03896v1 (CC BY 4.0), distributed by arXiv under the Creative Commons Attribution 4.0 International licence, http://creativecommons.org/licenses/by/4.0/, as recorded in the arXiv licence metadata for this exact version and shown on the abstract page https://arxiv.org/abs/2609.03896v1. Full licence notice: \texttt{licenses/arxiv-2609.03896-CC-BY-4.0.txt}.\par\smallskip

\noindent\textit{Editorial scope.} Counted unit: the article's theorem prop:clock-fta (source0.tex lines 567--570). The article's complete original proof of it is delivered with it (source0.tex lines 572--592, one \texttt{proof} environment of 21 lines). No mathematics is written, repaired or re-derived: the statement and the proof are the article's own lines, in the article's own notation, and no retained line is substituted or re-flowed.  Retained uncounted as the source-local context the proof needs: lines 512--520; lines 537--545.  Nothing was omitted from any retained span, so the package carries no omission record.  No retained line contains a citation, so the wrapper carries no bibliography and no bibliography entry.  No retained line contains a figure environment, an image-inclusion command or drawing code, so no image is delivered and the package declares no asset dependency.  Editorial formatting only, in addition: an AMS \texttt{amsart} preamble that restates the article's own declarations and macros used by the retained lines, namely the environments \texttt{lemma}, \texttt{proposition}, \texttt{theorem} on separate counters, together with the standard packages those lines need.  The retained environments print in this document as Lemma 1, Proposition 1, Theorem 1 in order of appearance, whereas the article prints the same items with the numbering of its own full document.  The complete native source of the article as distributed in the arXiv e-print for this exact version is bundled unchanged as \texttt{sources/209323/source0.tex}.
\begin{lemma}[Primewise additivity of the valuation readout]
\label{lem:valuation-additive}
For every prime $p$ and positive integers $a,b$,
\begin{equation}
V_p(ab)=V_p(a)+V_p(b).
\label{eq:valuation-additive}
\end{equation}
Hence $\V(ab)=\V(a)+\V(b)$ coordinatewise.
\end{lemma}

\begin{proposition}[Non-repetition of valuation rows]
\label{prop:valuation-nonrepetition}
For positive integers $n$ and $m$,
\begin{equation}
\V(n)=\V(m)\quad\Longrightarrow\quad n=m.
\label{eq:valuation-injective}
\end{equation}
Thus the valuation trajectory $\V(1),\V(2),\V(3),\ldots$ has no repeated row.
\end{proposition}

\begin{theorem}[Fundamental theorem of arithmetic, clock form]
\label{prop:clock-fta}
Every integer $n>1$ is a product of primes, and that product is unique up to the order of its factors.
\end{theorem}

\begin{proof}
For any finite-support nonnegative vector $\mathbf v=(v_p)_p$, form the integer
\[
F(\mathbf v):=\prod_p p^{v_p}.
\]
For each prime $p$, $\V(p)=\mathbf e_p$. Repeated application of Lemma~\ref{lem:valuation-additive} therefore gives
\begin{equation}
\V(F(\mathbf v))=\sum_p v_p\mathbf e_p=\mathbf v.
\label{eq:valuation-synthesis}
\end{equation}

For a fixed positive integer $n$, the row $\V(n)$ has finite support, since $V_p(n)=0$ whenever $p>n$. Taking $\mathbf v=\V(n)$ in~\eqref{eq:valuation-synthesis} gives
\[
\V(F(\V(n)))=\V(n).
\]
Both $F(\V(n))$ and $n$ are integer times on the same trajectory, so Proposition~\ref{prop:valuation-nonrepetition} forces
\[
n=F(\V(n))=\prod_p p^{V_p(n)}.
\]
This proves existence of a prime factorization directly from the valuation row. If also $n=\prod_p p^{a_p}$ for a finite-support family of nonnegative integers $(a_p)_p$, then additivity gives $\V(n)=(a_p)_p$. Hence $a_p=V_p(n)$ for every prime $p$, proving uniqueness up to the order of the factors.
\end{proof}

\end{document}
